% !TEX root = z-main.tex

Modelado 3D: fundamentos y aplicaciones

El modelado 3D se definió como una técnica de creación digital que permitió construir representaciones tridimensionales de objetos, espacios o personajes mediante herramientas computacionales. Fue considerado una forma de escultura digital, en la cual se partía de formas básicas —como cubos, planos o esferas— y se añadía geometría para generar volumen y detalle \cite{Teijeiro2023}. Este proceso combinó aspectos técnicos y artísticos, pues requirió tanto del dominio de programas especializados como de la comprensión de la forma, el espacio y los elementos visuales.

Las aplicaciones del modelado 3D abarcaron múltiples sectores, entre ellos el cine, los videojuegos, la arquitectura, el diseño industrial y, de forma creciente, la museología y la arqueología. En este último ámbito, los modelos tridimensionales facilitaron la reconstrucción de estructuras antiguas con precisión, favoreciendo su estudio y preservación, especialmente cuando los restos físicos se encontraron fragmentados, deteriorados o, como en el caso de Kaminaljuyú, permanecieron ocultos bajo la urbanización contemporánea. Dichas reconstrucciones se integraron además en entornos de realidad aumentada, ofreciendo experiencias inmersivas independientes del estado físico del sitio.

Realidad aumentada como herramienta educativa y cultural

La realidad aumentada (RA) se describió como una tecnología que permitió superponer elementos digitales —modelos tridimensionales, imágenes, textos o sonidos— sobre la percepción del mundo real en tiempo real. A diferencia de la realidad virtual, que construyó un entorno completamente simulado, la RA combinó el espacio físico con información digital interactiva, enriqueciendo la experiencia del usuario sin aislarlo de su contexto \cite{Torres2013}.

En el ámbito artístico-cultural, la RA abrió nuevas posibilidades para la presentación, interpretación y difusión del patrimonio. Museos y centros de interpretación la adoptaron como parte de sus estrategias museográficas, posibilitando el acceso a información ampliada sobre objetos culturales, la simulación de entornos históricos y la visualización de reconstrucciones virtuales de elementos desaparecidos o dañados. Estas aplicaciones fortalecieron la accesibilidad y la comprensión del contenido expuesto \cite{Torres2013}.

Desde la educación, la RA demostró ser una herramienta eficaz. Su incorporación en procesos de enseñanza-aprendizaje promovió la participación activa, el aprendizaje experiencial y la comprensión de conceptos complejos mediante la interacción directa con representaciones visuales y tridimensionales. Investigaciones recientes evidenciaron que la RA incrementó la motivación del estudiante, mejoró la retención del conocimiento y favoreció el aprendizaje significativo, especialmente en contextos donde los contenidos resultaron difíciles de representar de manera tradicional \cite{Barroso2022,Sevilla2017}.